\begin{afterword}
\summaryhead

The post takes its thesis from John Tooby and Leda Cosmides: ``Individual organisms are best thought of as adaptation-executers rather than as fitness-maximizers.'' Its example is taste. Taste buds that once led our ancestors to scarce sugar and fat now lead people in rich countries to ice cream and chocolate, which harm them. No one deliberately maximizing genetic fitness would eat a cookie unless starving.

To explain the distinction the post uses a screwdriver. Its cause is the designer's mind; its shape is the only one of these things found in the tool itself; its consequence is what it actually does; its meaning is what a user takes it to be for. The same four headings are then applied to the taste bud: its cause is a long line of ancestors who liked sugar and had children; its shape is a sensor linked to reinforcement; its consequence today can be overeating and fewer children; its meaning is up to the eater. The author rejects the looser formula that people do what would have spread their genes among hunter-gatherers, since no one asks that question, and ends by repeating the epigraph.

\responsehead

The post explains a sound distinction in plain terms and credits its slogan to Tooby and Cosmides. The four headings, applied first to a screwdriver and then to a taste bud, make the argument easy to follow, and its one figure, 2763 generations, matches the calculation in the post it links. Its one refinement, that people do not do ``what would have propagated our genes'' because they do not ask that question, is to the point.

What the reader is not told is where the distinction comes from in biology. Separating the history that explains a trait from the mechanism by which it works is what biologists call ultimate and proximate causation, terms Ernst Mayr popularized in 1961. Taste buds suited to scarce food, in a world of abundant food, are a textbook case of what is now called evolutionary mismatch. The post uses neither term. Nothing is lost for the argument, but a reader who wants to follow the idea further has no name to look for.

In short: a clear exposition of an established distinction, credited for its slogan, but without the biologists' names for the idea.
\end{afterword}
